
\section{Related Work}

\paragraph{Open-domain QA}
% odqa 설명
Open-domain QA is the task of accurately answering a query by sourcing for query-relevant documents, and then interpreting them to provide answers~\cite{DBLP:conf/acl/ChenFWB17, ODQA/survey}, which, thus, generally involves two modules: a retriever~\cite{DPR, ANCE} and a reader~\cite{BERTserini, FiD, DBLP:conf/emnlp/JeongBCHP23}. 
Along with the emergence of LLMs with superior reasoning capabilities thanks to their billion-sized parameters~\cite{EmergentLLMs}, a synergy between LLMs and retrievers has led to significant advancements~\cite{internet-QA, ram2023ralm}. Specifically, this integration has been shown to enhance Open-domain QA by mitigating the hallucination problem from LLMs through strengthened reasoning abilities of the reader, as well as utilizing the retrieved, external documents~\cite{DBLP:conf/emnlp/ChoSJP23}. Despite these advancements for single-hop retrieval-augmented LLMs, however, the complexity of some queries needs a more complex strategy.

\paragraph{Multi-hop QA}
Multi-hop QA is an extension of conventional Open-domain QA, which additionally requires the system to comprehensively gather and contextualize information from multiple documents (often iteratively), to answer more complex queries~\cite{musique, hotpotqa}. In the realm of multi-hop QA, the approach to iteratively access both LLMs and the retrieval module is generally employed. Specifically,~\citet{dsp},~\citet{self-ask},~\citet{DBLP:conf/ecir/PereiraFLN23} and~\citet{decomp} proposed to first decompose the multi-hop queries into simpler single-hop queries, repeatedly access the LLMs and retriever to solve these sub-queries, and merge their solutions to formulate a complete answer. In contrast to this decomposition-based approach, other recent studies, such as~\citet{react} and~\citet{ircot}, explored the interleaving of Chain-of-Thought reasoning~\cite{cot} — a method where a logical sequence of thoughts is generated — with document retrieval, repeatedly applying this process until the reasoning chain generates the answer. In addition,~\citet{ActiveRetrieval} introduced an approach to repeatedly retrieving new documents if the tokens within generated sentences have low confidence. However, the aforementioned methods overlooked the fact that, in real-world scenarios, queries are of a wide variety of complexities. Therefore, it would be largely inefficient to iteratively access LLMs and retrievers for every query, which might be simple enough with a single retrieval step or even only with an LLM itself. 

\paragraph{Adaptive Retrieval}
To handle queries of varying complexities, the adaptive retrieval strategy aims to dynamically decide whether to retrieve documents or not, based on each query's complexity.
In this vein,~\citet{adaptiveRetrieval} proposed to decide the query's complexity level based on the frequency of its entities and suggested using the retrieval modules only when the frequency falls below a certain threshold. However, this approach, focusing solely on the binary decision of whether to retrieve or not, may not be sufficient for more complex queries that require multiple reasoning steps.
Additionally, \citet{DBLP:conf/emnlp/0003LSM21} proposed an approach that performs a fixed set of operations (retrieving, reading, and reranking) multiple times until the answer is derived for the given query, which is built upon traditional BERT-like LMs. However, unlike our Adaptive-RAG which pre-determines the query complexity and adapts the operational behavior of any off-the-shelf LLMs accordingly, this approach applies the same fixed operations to every query regardless of its complexity but also necessitates additional specific training to LMs. Concurrent to our work,~\citet{self-rag} suggested training a sophisticated model to dynamically retrieve, critique, and generate the text. Nevertheless, we argue that all the aforementioned adaptive retrieval methods that rely on a single model might be suboptimal in handling a variety of queries of a range of different complexities since they tend to be either overly simple or complex for all the input queries, which demands a new approach that can select the most suitable strategy of retrieval-augmented LLMs tailored to the query complexity.
